\section{正规空间}\label{section:ch04-normal-spaces}

现在我们来更全面地研究一下满足正规性公理的空间。从某种意义上说，这个术语``正规''并不是很恰当，因为所谓正规空间并不如我们所想像的那样理想。另一方面，我们所熟知的许多空间都满足这个公理，就像我们下面所见到的那样。它之所以重要，是因为在假设了正规性的情况下，我们所能证明的很多结论都是拓扑学中很重要的结论。其中 Urysohn 度量化定理和 Tietze 扩张定理就是这样的两个结论，在本章稍后的几节中，我们将会处理这些问题。

我们先来证明三个定理，它们给出了判断空间是否正规的重要条件。

\begin{theorem}\label{theorem:ch04-second-countable-regular-normal}
每一个有可数基的正则空间是正规的。
\end{theorem}

\begin{proof}
设 $X$ 是有可数基 $\mathcal{B}$ 的一个正则空间，$A$ 和 $B$ 是 $X$ 的两个无交闭子集。$A$ 中的每一个点 $x$ 都存在一个邻域 $U$ 与 $B$ 无交。应用正则性，选取 $x$ 的一个邻域 $V$，使得 $\overline{V}$ 包含于 $U$。最后选取 $\mathcal{B}$ 中的一个包含着 $x$ 的元素，使得它包含于 $V$。通过对 $A$ 中每一个点 $x$ 选取基中的这样一个元素，我们就构造了 $A$ 的一个开覆盖，其中的每一个开集的闭包都与 $B$ 无交。因 $A$ 的这个覆盖是可数的，故我们可以选用正整数作为其下标，将其记为 $\{U_{n}\}$。

类似地，选取集合 $B$ 的一个可数开覆盖 $\{V_{n}\}$，使得每一个 $\overline{V}_{n}$ 都与 $A$ 无交。于是集合 $U=\bigcup U_{n}$ 和 $V=\bigcup V_{n}$ 就是分别包含 $A$ 和 $B$ 的开集，但它们未必是无交的。我们通过下面的方法来构造两个无交开集。给定 $n$，定义
\begin{align*}
U_{n}'=U_{n}-\bigcup_{i=1}^{n}\overline{V}_{i}\quad\text{和}\quad V_{n}'=V_{n}-\bigcup_{i=1}^{n}\overline{U}_{i}.
\end{align*}
注意，每一个集合 $U_{n}'$ 都是开集，因为它是一个开集 $U_{n}$ 和一个闭集 $\bigcup_{i=1}^{n}\overline{V}_{i}$ 的差。同样，每一个 $V_{n}'$ 也是开集。又因为 $A$ 中的每一个点 $x$ 都属于某一个 $U_{n}$，却不属于任何一个集合 $\overline{V}_{i}$，所以族 $\{U_{n}'\}$ 覆盖 $A$。同理，族 $\{V_{n}'\}$ 覆盖 $B$。如\cref{figure:ch04-normal-construction}所示。
\end{proof}

\par
\begin{figure}[htbp]
\centering
\includegraphics[width=0.7\textwidth]{ch04-fig32-1.png}
\caption{构造两个无交开集}\label{figure:ch04-normal-construction}
\end{figure}

最后，集合
\begin{align*}
U'=\bigcup_{n\in\mathbb{Z}_{+}}U_{n}'\quad\text{和}\quad V'=\bigcup_{n\in\mathbb{Z}_{+}}V_{n}'
\end{align*}
是无交的。因为若 $x\in U'\cap V'$，则对某一个 $j,k$，有 $x\in U_{j}'\cap V_{k}'$。不妨设 $j\leqslant k$。根据 $U_{j}'$ 的定义，有 $x\in U_{j}$。再根据 $j\leqslant k$ 以及 $V_{k}'$ 的定义，有 $x\notin\overline{U}_{j}$。对于 $j\geqslant k$，也有类似的矛盾产生。

\begin{theorem}\label{theorem:ch04-metrizable-normal}
每一个可度量化的空间是正规的。
\end{theorem}

\begin{proof}
设 $X$ 是一个度量空间，以 $d$ 为度量，$A$ 和 $B$ 是 $X$ 中的两个无交闭集。对于 $A$ 中的任意一个点 $a$，选取 $\varepsilon_{a}$ 使得球 $B(a,\varepsilon_{a})$ 与 $B$ 无交。类似地，对于 $B$ 中的任意一个点 $b$，取 $\varepsilon_{b}$ 使得球 $B(b,\varepsilon_{b})$ 与 $A$ 无交。定义
\begin{align*}
U=\bigcup_{a\in A}B(a,\varepsilon_{a}/2)\quad\text{和}\quad V=\bigcup_{b\in B}B(b,\varepsilon_{b}/2),
\end{align*}
于是 $U$ 和 $V$ 就是分别包含集合 $A$ 和 $B$ 的开集。我们断言：它们是无交的。因为若 $z\in U\cap V$，则存在 $a\in A$ 和 $b\in B$，使得
\begin{align*}
z\in B(a,\varepsilon_{a}/2)\cap B(b,\varepsilon_{b}/2).
\end{align*}
根据三角不等式可见 $d(a,b)<(\varepsilon_{a}+\varepsilon_{b})/2$。若 $\varepsilon_{a}\leqslant\varepsilon_{b}$，则 $d(a,b)<\varepsilon_{b}$，从而球 $B(b,\varepsilon_{b})$ 包含点 $a$。若 $\varepsilon_{a}>\varepsilon_{b}$，则 $d(a,b)<\varepsilon_{a}$，从而球 $B(a,\varepsilon_{a})$ 包含点 $b$。但这两种情形都是不可能的。
\end{proof}

\begin{theorem}\label{theorem:ch04-compact-hausdorff-normal}
每一个紧致的 Hausdorff 空间都是正规的。
\end{theorem}

\begin{proof}
设 $X$ 是一个紧致的 Hausdorff 空间。实际上我们已经证明了 $X$ 是正则的。因为若 $x$ 是 $X$ 的一个点，$B$ 是 $X$ 中不包含 $x$ 的一个闭集，则 $B$ 是紧致的。于是应用第 3 章引理 26.4，就证明了存在无交的开集分别包含 $x$ 和 $B$。

采用与上述引理中基本上相同的论证，就可以证得 $X$ 是正规的：给定 $X$ 中无交的闭集 $A$ 和 $B$，对于 $A$ 中的每一点 $a$，分别选取包含 $a$ 和 $B$ 的无交开集 $U_{a}$ 和 $V_{a}$。（这里用到了 $X$ 的正则性。）于是集簇 $\{U_{a}\}$ 覆盖 $A$。再根据 $A$ 的紧致性可见，可以从中选出有限个集合 $U_{a_{1}},\cdots,U_{a_{m}}$ 覆盖 $A$。从而
\begin{align*}
U=U_{a_{1}}\cup\cdots\cup U_{a_{m}}\quad\text{和}\quad V=V_{a_{1}}\cap\cdots\cap V_{a_{m}}
\end{align*}
就是分别包含 $A$ 和 $B$ 的无交的开集。
\end{proof}

下面这个关于正规性的定理，在处理一些例子时是很有用的。

\begin{theorem}\label{theorem:ch04-well-ordered-normal}
每一个良序集 $X$ 在序拓扑下都是正规的。
\end{theorem}

事实上，每一个序拓扑都是正规的（见 [S-S] 的例 39），但是我们用不着这个更强的结果。

\begin{proof}
设 $X$ 是一个良序集，我们断言：每一个形如 $(x,y]$ 的区间都是 $X$ 中的开集：若 $X$ 有最大元 $y$，则 $(x,y]$ 恰好就是包含 $y$ 的一个基元素。若 $y$ 不是 $X$ 的最大元，则 $(x,y]$ 就等于开集 $(x,y')$，其中 $y'$ 为 $y$ 的直接后继元。

现在设 $A$ 和 $B$ 为 $X$ 中无交的两个闭集，并且 $A$ 和 $B$ 都不包含 $X$ 的最小元 $a_{0}$。对于每一个 $a\in A$，总存在包含 $a$ 但与 $B$ 无交的基元素，它包含某一个形如 $(x,a]$ 的区间（这里我们利用了 $a$ 不是 $X$ 的最小元这个事实）。这样对于每一个 $a\in A$，便有一个与 $B$ 无交的区间 $(x_{a},a]$。类似地，对于每一个 $b\in B$，也有一个与 $A$ 无交的区间 $(y_{b},b]$。于是，集合
\begin{align*}
U=\bigcup_{a\in A}(x_{a},a]\quad\text{和}\quad V=\bigcup_{b\in B}(y_{b},b]
\end{align*}
便是分别包含 $A$ 和 $B$ 的开集，下面我们证明它们无交。因为若 $z\in U\cap V$，则对于某一个 $a\in A$ 和 $b\in B$，有 $z\in(x_{a},a]\cap(y_{b},b]$。不妨假定 $a<b$，若 $a\leqslant y_{b}$，则两个区间无交。若 $a>y_{b}$，则 $a\in(y_{b},b]$。这和 $(y_{b},b]$ 与 $A$ 无交矛盾。同理，对于 $a>b$，同样可以得出矛盾。

最后，设 $A$ 和 $B$ 是 $X$ 中无交的两个闭集，并且 $A$ 包含 $X$ 的最小元 $a_{0}$。集合 $\{a_{0}\}$ 在 $X$ 中是既开又闭的。由上一段的结果知，存在分别包含闭集 $A-\{a_{0}\}$ 和 $B$ 的无交开集 $U$ 和 $V$，于是 $U\cup\{a_{0}\}$ 和 $V$ 就是分别包含 $A$ 和 $B$ 的无交的开集。
\end{proof}

\begin{example}\label{example:ch04-rj-not-normal}
若 $J$ 是不可数的，则积空间 $\mathbb{R}^{J}$ 就不是正规的。这个结论的证明略有困难，我们将它留作一道具有挑战性的习题（见习题 9）。

举这个例子有三个目的。它指出了正则空间 $\mathbb{R}^{J}$ 未必是正规的。它说明了正规空间的子空间未必是正规的，因为 $\mathbb{R}^{J}$ 同胚于 $[0,1]^{J}$ 的子空间 $(0,1)^{J}$，而 $[0,1]^{J}$（根据 Tychonoff 定理）是紧致的 Hausdorff 空间，因而是正规的。最后，它还说明了正规空间的不可数积未必是正规的。于是，剩下的问题只有正规空间的有限积或可数积是否正规了。
\end{example}

\begin{example}\label{example:ch04-somega-product-not-normal}
积空间 $S_{\Omega}\times\overline{S}_{\Omega}$ 不是正规的。

考虑具有序拓扑的良序集 $\overline{S}_{\Omega}$，及其具有子空间拓扑（与序拓扑相同）的子集 $S_{\Omega}$。根据\cref{theorem:ch04-well-ordered-normal}，这两个空间都是正规的。我们将证明积空间 $S_{\Omega}\times\overline{S}_{\Omega}$ 不是正规的。

这个例子可以达到三个目的：第一，它表明正则空间未必是正规的。因为 $S_{\Omega}\times\overline{S}_{\Omega}$ 是正则空间的积，因而是正则的。第二，它表明正规空间的子空间未必是正规的。因为 $S_{\Omega}\times\overline{S}_{\Omega}$ 是 $\overline{S}_{\Omega}\times\overline{S}_{\Omega}$ 的一个子空间，而 $\overline{S}_{\Omega}\times\overline{S}_{\Omega}$ 是紧致的 Hausdorff 空间，因而是正规的。第三，它表明两个正规空间的积空间未必是正规的。

首先考虑空间 $\overline{S}_{\Omega}\times\overline{S}_{\Omega}$ 和它的``对角线''$\Delta=\{\,x\times x\mid x\in\overline{S}_{\Omega}\,\}$。因为 $\overline{S}_{\Omega}$ 是 Hausdorff 的，所以 $\Delta$ 是 $\overline{S}_{\Omega}\times\overline{S}_{\Omega}$ 的一个闭集：这是因为若 $U$ 和 $V$ 分别是 $x$ 和 $y$ 的无交的邻域，则 $U\times V$ 是 $x\times y$ 的一个邻域，并且与 $\Delta$ 无交。

因此，在子空间 $S_{\Omega}\times\overline{S}_{\Omega}$ 中，集合
\begin{align*}
A=\Delta\cap(S_{\Omega}\times\overline{S}_{\Omega})=\Delta-\{\Omega\times\Omega\}
\end{align*}
是闭的。同样，集合
\begin{align*}
B=S_{\Omega}\times\{\Omega\}
\end{align*}
是积空间中的一个``薄片''，所以在 $S_{\Omega}\times\overline{S}_{\Omega}$ 中是闭的。并且集合 $A$ 和 $B$ 是无交的。我们断言：若在 $S_{\Omega}\times\overline{S}_{\Omega}$ 中存在分别包含 $A$ 和 $B$ 的无交的开集 $U$ 和 $V$，则必导致矛盾出现。如\cref{figure:ch04-somega-product}所示。
\end{example}

\par
\begin{figure}[htbp]
\centering
\includegraphics[width=0.5\textwidth]{ch04-fig32-2.png}
\caption{$S_{\Omega}\times\overline{S}_{\Omega}$ 不是正规的}\label{figure:ch04-somega-product}
\end{figure}

给定 $x\in S_{\Omega}$，考虑竖直的一个``薄片'' $x\times\overline{S}_{\Omega}$。我们证明存在某一点 $\beta$，$x<\beta<\Omega$，使得 $x\times\beta$ 落在 $U$ 的外面。因为如果 $U$ 包含所有的点 $x\times\beta(x<\beta<\Omega)$，那么薄片的顶点 $x\times\Omega$ 就是 $U$ 的一个极限点。但这是不可能的。因为 $V$ 是与 $U$ 无交的开集，并且包含这个顶点。

选取 $\beta(x)$ 为 $S_{\Omega}$ 中使得 $x<\beta(x)<\Omega$ 并且使得 $x\times\beta(x)$ 落在 $U$ 外面的最小元素。定义 $S_{\Omega}$ 中点的一个序列如下：设 $x_{1}$ 为 $S_{\Omega}$ 中任意一点，令 $x_{2}=\beta(x_{1})$。对于一般情形，$x_{n+1}=\beta(x_{n})$。我们有
\begin{align*}
x_{1}<x_{2}<\cdots.
\end{align*}
因为对于所有的 $x$，$\beta(x)>x$。集合 $\{x_{n}\}$ 是可数的，因而在 $S_{\Omega}$ 中有一个上界。设 $b\in S_{\Omega}$ 是它的上确界。因为序列是递增的，那么它必收敛于它的上确界。即 $x_{n}\to b$。但是 $\beta(x_{n})=x_{n+1}$，于是 $\beta(x_{n})$ 也收敛于 $b$。因此，在积空间中
\begin{align*}
x_{n}\times\beta(x_{n})\to b\times b.
\end{align*}
参见\cref{figure:ch04-somega-sequence}。这就产生了矛盾。因为点 $b\times b$ 属于集合 $A$，并且集合 $A$ 包含在开集 $U$ 中，但是 $x_{n}\times\beta(x_{n})$ 中没有点在 $U$ 中。

\par
\begin{figure}[htbp]
\centering
\includegraphics[width=0.5\textwidth]{ch04-fig32-3.png}
\caption{$x_{n}\times\beta(x_{n})\to b\times b$}\label{figure:ch04-somega-sequence}
\end{figure}

\subsection*{习题}

\begin{enumerate}
\item 证明正规空间的闭子空间是正规的。
\item 证明：若 $\prod X_{\alpha}$ 是 Hausdorff 的，或者正则的，或者正规的，则 $X_{\alpha}$ 也相应地是 Hausdorff 的，或者正则的，或者正规的。（假设每一个 $X_{\alpha}$ 都是非空的。）
\item 证明每一个局部紧致的 Hausdorff 空间是正则的。
\item 证明每一个正则的 Lindelöf 空间是正规的。
\item 在积拓扑下，$\mathbb{R}^{\omega}$ 是正规的吗？在一致拓扑下呢？

$\mathbb{R}^{\omega}$ 在箱拓扑下是否是正规的，尚属未知。Mary-Ellen Rudin 已证明了在连续统假设的情况下，答案是肯定的 [RM]。事实上，她还证明了这个空间满足一个称之为仿紧致性的更强的性质。
\item 空间 $X$ 称为\textbf{完全正规的}（completely normal），如果 $X$ 的每一个子空间都是正规的。证明：$X$ 是完全正规的当且仅当 $X$ 中的每一对分离集 $A,B$（即 $A$ 和 $B$ 满足：$\overline{A}\cap B=\varnothing$ 和 $A\cap\overline{B}=\varnothing$），存在无交的开集分别包含着它们。[提示：若 $X$ 是完全正规的，考虑 $X-(\overline{A}\cap\overline{B})$。]
\item 以下空间中的哪一个是完全正规的？证明你的结论。
\begin{enumerate}
\item 完全正规空间的子空间.
\item 两个完全正规空间的积空间.
\item 序拓扑下的良序集.
\item 可度量化空间.
\item 紧致的 Hausdorff 空间.
\item 具有可数基的正则空间.
\item 空间 $\mathbb{R}_{\ell}$.
\end{enumerate}
\item 证明以下结论：

\textbf{定理}\quad 每一个线性连续统 $X$ 是正规的.
\begin{enumerate}
\item 设 $C$ 是 $X$ 的一个非空闭子集。如果 $U$ 是 $X-C$ 的一个分支，证明 $U$ 是一个形如 $(c,c')$，$(c,\infty)$ 或者 $(-\infty,c)$ 的集合，其中 $c,c'\in C$.
\item 设 $A$ 和 $B$ 是 $X$ 的两个无交闭子集。对于 $X-A\cup B$ 的那些一个端点在 $A$ 中，另一个端点在 $B$ 中的分支 $W$，选取 $W$ 的一个点 $c_{W}$。证明这些点 $c_{W}$ 构成的集合 $C$ 是闭的.
\item 证明：若 $V$ 是 $X-C$ 的一个分支，则 $V$ 与 $A$ 和 $B$ 都无交.
\end{enumerate}
\item 证明以下结论：

\textbf{定理}\quad 若 $J$ 是不可数的，则 $\mathbb{R}^{J}$ 不是正规的.

证明（这个证明由 A. H. Stone 给出，[S-S] 中作了改写）：设 $X=(\mathbb{Z}_{+})^{J}$。只要证明 $X$ 不是正规的就可以了，因为 $X$ 是 $\mathbb{R}^{J}$ 的一个闭子集。对于 $X$ 的元素，我们使用函数记号，于是 $X$ 中的元素就是一个函数 $\boldsymbol{x}\colon J\to\mathbb{Z}_{+}$。
\begin{enumerate}
\item 设 $\boldsymbol{x}\in X$，$B$ 是 $J$ 的有限子集。用 $U(\boldsymbol{x},B)$ 表示 $X$ 中所有满足 $\boldsymbol{y}(\alpha)=\boldsymbol{x}(\alpha)$ 的点 $\boldsymbol{y}$ 构成的集合，其中 $\alpha\in B$。证明：所有的集合 $U(\boldsymbol{x},B)$ 构成了 $X$ 的一个基.
\item 定义 $P_{n}$ 为 $X$ 中所有那些点 $\boldsymbol{x}$ 构成的集合，这些点 $\boldsymbol{x}$ 在集合 $J-\boldsymbol{x}^{-1}(n)$ 上是单射。证明：$P_{1}$ 和 $P_{2}$ 都是闭集，并且无交.
\item 设 $U$ 和 $V$ 是分别包含 $P_{1}$ 和 $P_{2}$ 的开集。给定 $J$ 的子序列 $\alpha_{1},\alpha_{2},\cdots$（其元素两两不同）和一个整数序列
\begin{align*}
0=n_{0}<n_{1}<n_{2}<\cdots,
\end{align*}
对于每一个 $i\geqslant 1$，令
\begin{align*}
B_{i}=\{\alpha_{1},\cdots,\alpha_{n_{i}}\}
\end{align*}
并且用等式
\begin{align*}
\boldsymbol{x}_{i}(\alpha_{j})=j\quad\text{如果 }1\leqslant j\leqslant n_{i-1},
\end{align*}
\begin{align*}
\boldsymbol{x}_{i}(\alpha)=1\quad\text{当 }\alpha\text{ 取其他值时}.
\end{align*}
定义 $\boldsymbol{x}_{i}\in X$。证明：可以选取序列 $\alpha_{j}$ 和 $n_{j}$，使得对于每一个 $i$ 以下包含关系成立
\begin{align*}
U(\boldsymbol{x}_{i},B_{i})\subset U.
\end{align*}
[提示：注意，对于所有的 $\alpha$ 有 $\boldsymbol{x}_{1}(\alpha)=1$。选取 $B_{1}$ 使得 $U(\boldsymbol{x}_{1},B_{1})\subset U$。]
\item 设 $A$ 为上步 (c) 中的集合 $\{\alpha_{1},\alpha_{2},\cdots\}$。定义 $\boldsymbol{y}\colon J\to\mathbb{Z}_{+}$ 为
\begin{align*}
\boldsymbol{y}(\alpha_{j})=j\quad\text{当 }\alpha_{j}\in A\text{ 时},
\end{align*}
\begin{align*}
\boldsymbol{y}(\alpha)=2\quad\text{当 }\alpha\text{ 取其他值时}.
\end{align*}
选取 $B$ 使得 $U(\boldsymbol{y},B)\subset V$。再取 $i$ 使得 $B\cap A$ 包含在集合 $B_{i}$ 中。证明
\begin{align*}
U(\boldsymbol{x}_{i+1},B_{i+1})\cap U(\boldsymbol{y},B)
\end{align*}
非空.
\end{enumerate}
\item 每一个拓扑群都是正规的吗？
\end{enumerate}